\documentclass{article}

\usepackage{tabularx}
\usepackage{booktabs}

\title{Problem Statement and Goals\\\progname}

\author{\authname}

\date{}

\input{../Comments.text}
\input{../Common.text}

\begin{document}

\maketitle

\begin{table}[hp]
\caption{Revision History} \label{TblRevisionHistory}
\begin{tabularx}{\textwidth}{llX}
\toprule
\textbf{Date} & \textbf{Developer(s)} & \textbf{Change}\\
\midrule
September 25, 2026 & Ethan Walsh \& Jinwoo Hong & Creation of initial document \\
\bottomrule
\end{tabularx}
\end{table}

\section{Problem Statement}

\subsection{Problem}

As the use of artificial intelligence inference becomes more common, the most 
popular option is to use a commercial large language model from a major inference 
provider via a web interface or API (application programming interface). This approach is subject to variable costs 
and insufficient data privacy protections. Although some users have opted for 
open-source models on local devices, these models are often limited in quality 
due to hardware restrictions. Users currently cannot access the benefits of 
both approaches: low-cost, high-quality, private artificial intelligence inference. 
Addressing this gap would create more democratized access to inference.

\subsection{Inputs and Outputs}

Inputs:

\begin{itemize}
    \item A large language model
    \item A user's local device(s)
    \item User inputs to the interface (clicks)
\end{itemize}

Outputs:
\begin{itemize}
    \item User Interface (more specifically, output from the model)
\end{itemize}

\subsection{Stakeholders}

Stakeholders include, but are not limited to:

\begin{itemize}
    \item Local AI users who run models on their own device(s)
    \item Small companies in regulated industries (more specifically, industries that place high restrictions on data privacy)
    \item Low-budget and smaller AI labs focusing on research related to AI models
\end{itemize}

An example of a small company in a regulated industry would be a local medical clinic. This clinic may require AI or computer vision based analysis of patient scan results but this data must remain private.

\subsection{Environment}

\subsubsection{Hardware}

Hardware used for this project includes consumer devices such as personal computers and laptops. Mobile phones may also be used in the final iteration of the project but is not in-scope for the initial proof of concept.
Hardware constraints have yet to be specified due to the project being in its early stages. As design and development progresses, hardware constraints will be determined by the team and listed.

\subsubsection{Software}

Software must be compatible with various operating systems to be compatible on consumer devices. This includes Windows, Linux, and MacOS.

\section{Goals}

\subsection{Improved Local Performance Outcomes:}
The system's model will achieve better performance than if the equivalent model is only run on one device. The aim of this goal is to generally improve local AI outcomes and performance so that the product can become a viable alternative to commercial LLMs for simple tasks. 'Better' will be measured in a variety of metrics including:
\begin{itemize}
    \item Response time
    \item Throughput (tokens per second)
    \item Model benchmarking
    \item Model size that can be accommodated (number of parameters and quantization)
\end{itemize}

\subsection{Model Parallelization}
This goal is the motivating question that brought about this project: can a model be parallelized across multiple local devices and run successfully by producing coherent output? This goal is defined as the system parallelizing an artificial intelligence model (more specifically, a large language model based on the transformer architecture) across more than one local device. Currently, this goal will be considered met when the system is able to split up the layers of a model across the pool of available devices.
    
\subsection{Cross-Platform Optimization}
The system is able to operate across different operating systems and platforms (e.g. Windows, Mac, and Linux) so that users may easily connect their devices. Achieving this will result in the user base not being limited based on their existing device type.

\subsection{Dynamic Allocation}
In the case goal 2.2 can be achieved (model parallelization), the system can dynamically allocate the parallelized model across the pool of available devices.
This is a key shortcoming of existing local distributed inference providers that only do one allocation pass at the start of the session. By achieving this goal, the system would be more efficient and resilient since it can continue to operate as devices drop or are added to the network.

\end{itemize}

\section{Stretch Goals}

\begin{itemize}
    \item The system can operate via both wired and wireless connection (v0: wired, v1+: wireless)
\end{itemize}

\wss{Goals you hope to get to, but not necessary.  They are also helpful if the
scope of the original project needs to be expanded.}

\section{Custom Documents (CDs)}

\begin{itemize}
    \item \textbf{Fall:} Design Thinking Report
    \item \textbf{Winter:} Performance Report
\end{itemize}

\section{Generative AI Use Disclosure}
Anthropic's Claude Sonnet 5 was used as a method to provide feedback on early revisions of the problem statement (rubric and example problem statements from the capstone GitLab were provided as context). No text generated by the model was used and only general suggestions were used as guidelines for improvement.
Google's AI (built into search functionality) was also used as a method to provide feedback on the environment section (however feedback was largely ignored and was deemed to be unhelpful).

\newpage{}

\section*{Appendix --- Reflection}

\input{../Reflection.text}

\begin{enumerate}
    \item What went well while writing this deliverable? 
    \item What pain points did you experience during this deliverable, and how
    did you resolve them?
    \item How did you and your team adjust the scope of your goals to ensure
    they are suitable for a Capstone project (not overly ambitious but also of
    appropriate complexity for a senior design project)?
\end{enumerate} 

\end{document}